\documentclass{article}
\usepackage{graphicx}
\usepackage{hyperref}
\usepackage{booktabs}
\usepackage{geometry}
\geometry{a4paper, margin=1in}

\title{Foundations for Sentiment-Aware Trading: \\ Data Pipeline \& Baseline Validation}
\author{Abdurrahim Bayraktar}
\date{\today}

\begin{document}

\maketitle

\begin{abstract}
This report details the construction of a robust, scalable data pipeline and the establishment of a verified price-only baseline model for stock price prediction. The primary objective is to create a stable foundation for future sentiment-aware trading models. We systematically scaled a pooled LSTM model from single-stock analysis to over 500 stocks. During this process, we identified and corrected a critical data leakage issue (v7 to v10) that initially inflated accuracy to 81\%, establishing the true baseline at approximately 38\% accuracy for a 3-class prediction task. Furthermore, we present preliminary sentiment analysis results using FinBERT, confirming a significant same-day correlation between news sentiment and price returns, validating the potential for sentiment-based alpha.
\end{abstract}

\section{Introduction}
The ultimate goal of this project is to develop a stock trading system that effectively leverages financial news sentiment to predict price movements. However, integrating sentiment data into a predictive model requires a rigorous and reproducible foundation. This report focuses on the "pre-sentiment" phase: building a reliable data pipeline and establishing a trustworthy price-only baseline. Without a solid baseline, improvements from sentiment features cannot be accurately measured.

\section{The Data Pipeline: Efficiency \& Reproducibility}
To handle the scale of analyzing hundreds of stocks and thousands of news articles, we built a robust, cached data pipeline. This infrastructure is critical for rapid iteration and reproducibility.

\subsection{Architecture Overview}
The system is designed to minimize duplicate work and external API calls.
\begin{enumerate}
    \item \textbf{Data Fetching}: We utilize `yfinance` to retrieve historical price data and the FNSPID dataset for financial news.
    \item \textbf{Price Cache}: Raw price data is stored in a structured pickle cache (`PriceCache`). This prevents rate-limiting issues from external providers and ensures all experiments run on the exact same historical snapshot.
    \item \textbf{Sentiment Cache}: Sentiment extraction is computationally expensive. We cache the output of the FinBERT model (sentiment scores and labels) so that once a news item is processed, it never needs to be re-computed. This is crucial for iterating on downstream models without re-running the GPU-intensive inference.
    \item \textbf{Feature Cache}: Technical indicators (over 100 features) are pre-calculated and cached (`ComprehensiveIndicators`). This reduced experiment startup time from hours to minutes.
\end{enumerate}

\subsection{Strict Validation Methodology}
A key innovation in our pipeline is the implementation of \textbf{Temporal Splitting} at the individual stock level before pooling. 
\begin{itemize}
    \item \textbf{Train}: Data prior to 2022-01-01.
    \item \textbf{Validation}: Data from 2022-01-01 to 2023-06-01.
    \item \textbf{Test}: Data from 2023-06-01 to 2024-12-31.
\end{itemize}
This ensures that the model never trains on future data, a common pitfall in financial ML (look-ahead bias).

\section{Baseline Experiments: The Quest for a Stable Reference}
We conducted a series of experiments (v1 through v10) to establish a price-only baseline. This journey highlighted the extreme sensitivity of financial models to data leakage.

\subsection{Scaling Up (v1 - v6)}
We began with a single stock (MSFT) which yielded random results (approx. 33\% accuracy). By scaling to a "pooled" model architecture—training a single LSTM on data from over 500 stocks simultaneously—we successfully forced the model to learn generalizable patterns rather than memorizing a specific ticker's history. This approach raised our test accuracy to a respectable \textbf{51.57\%} (v6).

\subsection{The Leakage Lesson (v7)}
In iteration v7, we introduced "Domain Features" including Ichimoku Cloud indicators. This iteration achieved a stunning \textbf{81.45\%} test accuracy. 
Upon investigation, this was identified as \textbf{Data Leakage}. The feature `ichimoku\_ICS\_26` (Chikou Span) plots the current closing price 26 periods backwards. By including this feature, we inadvertently gave the model access to the future price 26 days ahead.

\subsection{Verified Baseline (v10)}
Iteration v10 rectified this error by removing the leaking features and enforcing strict temporal boundaries. 
\begin{itemize}
    \item \textbf{Model}: Attention-based LSTM.
    \item \textbf{Input}: 108 Technical Indicators (Trend, Momentum, Volatility, Volume).
    \item \textbf{Result}: \textbf{38.39\% Test Accuracy} (3-class classification: Up/Down/Neutral).
\end{itemize}
While significantly lower than the v7 "miracle," v10 represents the \textit{truth}. It performs better than a random guess (33.3\%) and the "Zero Rule" baseline (always predicting the majority class), confirming that technical analysis alone contains weak signal. This is the verified control group against which our sentiment models will compete.

\section{Sentiment Feasibility Study}
Parallel to the baseline work, we validated the quality of our sentiment data using the FinBERT model on the top 40 most covered stocks.

\subsection{Methodology}
We extracted sentiment from headlines and computed the correlation between the daily aggregated sentiment score and stock returns.
\begin{itemize}
    \item \textbf{Same-Day Correlation}: $Corr(Sentiment_t, Return_t)$
    \item \textbf{Next-Day Correlation}: $Corr(Sentiment_t, Return_{t+1})$
\end{itemize}

\subsection{Results}
\begin{table}[h]
\centering
\begin{tabular}{lcc}
\toprule
Ticker & Same-Day Correlation & Next-Day Correlation \\
\midrule
MSFT & \textbf{0.347} & -0.003 \\
DIS & 0.087 & -0.026 \\
WMT & 0.079 & -0.040 \\
\bottomrule
\end{tabular}
\caption{Correlation Analysis for Key Tickers}
\end{table}

\textbf{Key Finding}: There is a strong, positive correlation (0.347 for MSFT) on the \textit{same day}. This confirms that news moves markets and our sentiment extraction is accurate. However, the near-zero correlation for the \textit{next day} indicates that the market is efficient—news is priced in almost instantly.

\section{Conclusion \& Future Work}
We have successfully built a high-performance data pipeline and established a rigorously verified baseline of ~38\% accuracy. Taking into account the efficiency of the market (demonstrated by the lack of simple next-day correlation), our future work will focus on:
\begin{enumerate}
    \item \textbf{Intraday Analysis}: Attempting to capture the market reaction within minutes or hours of a news release, rather than waiting for the daily close.
    \item \textbf{Regime Shifts}: Using sentiment spread to predict volatility regimes rather than directional price.
    \item \textbf{Advanced Architectures}: Utilizing the clean V10 dataset to train Transformer-based models that might capture non-linear relationships better than the current LSTM.
\end{enumerate}

\end{document}
